% TOP_PROBLEM: {"id": 20002510, "title": "Normalization-aware leading quintic amplitudes at the Gepner point", "category_id": 16, "category": {"id": 16, "name": "physics", "display_name": "Mathematical Physics", "description": "Problems at the intersection of mathematics and physics.", "slug": "physics", "order_index": 16, "created_at": "2026-07-31T15:26:25.671Z"}, "status": "partially_solved", "statusQualification": "The source does not fix the flat coordinate, period normalization or holomorphic ambiguity. These choices must be recorded for numerical coefficients to be meaningful; z=0 is retained in the displayed family convention. The catalogue title/status was paired with a machine-generated research attempt; that attempt is not peer-reviewed evidence of settlement. The statement covers the original source questions rather than only the title’s subcase."}
% FIRST_POSTED: 2026-10-11
% ENTRY_KIND: statement_only
% SUBMISSION: {"kind": "statement_only", "category": "top_problems", "problem_id": "20002510", "proposal_type": "established", "source": "American Institute of Mathematics, Gromov-Witten invariants and number theory, source item 4.1. http://aimpl.org/gromwitnumthry/4/", "status_evidence": "Inherited source-catalogue status partially_solved preserved. Background has machine-generated research claims explicitly marked not peer reviewed; these are not independent status evidence. Primary source accessed 11 October 2026; compared against the supplied catalogue excerpt. Historical wording and partial-result remarks were checked for scope; no comprehensive current-literature or proof-verification audit was performed.", "status_checked_on": "2026-10-11", "duplicate_check": "Full 15,773-record source registry checked for identical normalized statements; available site registry checked for identical normalized title/statement and canonical-ID collision. Canonical ID is absent from the site registry. No mathematical-equivalence literature audit performed. Exact statement matches: 20000721.", "attribution": "Proposer not named in the AIM problem entry. Statement preparation: OpenAI GPT-6 Codex assistant, generation 11 October 2026; source-matching and notation editing by the assistant, no independent human mathematical verification. AIM Problem Lists source states Creative Commons Attribution–ShareAlike; preserve its attribution and applicable licence."}
% SOURCE_STATEMENT: "Can we calculate the leading terms of $F_g^B(q)$ at the orbifold point $z=0$ of the quintic\n\\[\n(\\sum x_i^5 + z \\prod x_i =0 )/ \\mathbb{Z}_5^3?\n\\]"
% Preparation and final self-audit: CONTRIBUTING.md prompts C and D; source/audit artifacts retained in the local physics_statements_2026_10_11 batch.
% Imported from: physics_statements.tex; UnsolvedMath ID 20002510
% !TeX program = lualatex
\documentclass[11pt]{article}
\usepackage[margin=25mm]{geometry}
\usepackage{fontspec}
\setmainfont{Latin Modern Roman}
\usepackage{amsmath,amssymb,mathtools}
\usepackage{unicode-math}
\setmathfont{Latin Modern Math}
\usepackage{xurl,hyperref}
\hypersetup{colorlinks=true,urlcolor=blue}
\setcounter{secnumdepth}{0}
\setlength{\parindent}{0pt}
\setlength{\parskip}{5pt}
\setlength{\emergencystretch}{3em}
\newcommand{\sref}[2]{\hyperlink{source:#1:#2}{[S#2]}}
\newcommand{\cataloguescope}[1]{\paragraph{Recorded target.}#1}
\newcommand{\R}{\mathbb R}
\newcommand{\C}{\mathbb C}
\newcommand{\N}{\mathbb N}
\newcommand{\Z}{\mathbb Z}
\newcommand{\Q}{\mathbb Q}
\newcommand{\F}{\mathbb F}
\newcommand{\eps}{\varepsilon}
\begin{document}
% UnsolvedMath ID 20002510; source catalogue code AIM-PHYSICS-0018
\setcounter{section}{20002509}
\section{Normalization-aware leading quintic amplitudes at the Gepner point}\label{20002510}
{\small UnsolvedMath ID 20002510 · Mathematical Physics}
\subsection{Definitions and mathematical statement}

\paragraph{Shared notation.}$\mathbb N=\{1,2,\ldots\}$; the standard number fields have their usual meanings. All additional conventions are stated locally.


\paragraph{Definitions and notation.}The displayed family is the mirror-quintic model obtained from $\sum_{i=1}^5x_i^5+z\prod_{i=1}^5x_i=0$ by the indicated $(\mathbb Z/5\mathbb Z)^3$ quotient, with the geometric resolution/convention required for mirror calculations. $F_g^B$ denotes a genus-$g$ B-model amplitude. The variable $q$ is the source’s amplitude coordinate; its relation to $z$ and period normalization matters at the orbifold/Gepner point.

\paragraph{Catalogue wording retained for traceability.}
Can we calculate the leading terms of $F_g^B(q)$ at the orbifold point $z=0$ of the quintic
\[
(\sum x_i^5 + z \prod x_i =0 )/ \mathbb{Z}_5^3?
\]
\paragraph{Scope and source qualification.}The source does not fix the flat coordinate, period normalization or holomorphic ambiguity. These choices must be recorded for numerical coefficients to be meaningful; z=0 is retained in the displayed family convention. The catalogue title/status was paired with a machine-generated research attempt; that attempt is not peer-reviewed evidence of settlement. The statement covers the original source questions rather than only the title’s subcase.
\subsection{Short English statement}
Compute the leading B-model genus amplitudes of the mirror quintic near its orbifold point.
\subsection{Sources}
\begin{itemize}
\item[\hypertarget{source:20002510:1}{S1}] ULAM.AI and the UnsolvedMath Contributors, UnsolvedMath, record 20002510 (AIM-PHYSICS-0018). Catalogue curation under CC BY 4.0; definitions and English exposition added locally.
\url{https://huggingface.co/datasets/ulamai/UnsolvedMath}.
\item[\hypertarget{source:20002510:2}{S2}] American Institute of Mathematics, Gromov-Witten invariants and number theory, source item 4.1.
\url{http://aimpl.org/gromwitnumthry/4/}.
\end{itemize}
\label{end:20002510}
\end{document}
